% Anonymous two-page English CV template
% Compile with XeLaTeX and replace all example fields.
\documentclass[a4paper,11pt]{article}
\usepackage[a4paper,top=14mm,bottom=14mm,left=17mm,right=17mm]{geometry}
\usepackage{fontspec}
\usepackage{xcolor}
\usepackage{titlesec}
\usepackage{enumitem}
\usepackage[hidelinks]{hyperref}
\usepackage{tabularx}
\usepackage{fancyhdr}
\usepackage{ragged2e}
\setmainfont{Calibri}
\XeTeXgenerateactualtext=1
\definecolor{accent}{HTML}{3F489D}
\definecolor{secondary}{HTML}{7281C2}
\definecolor{panelbg}{HTML}{F0F2F8}
\definecolor{textgrey}{HTML}{24282B}
\definecolor{muted}{HTML}{5F686E}
\hypersetup{pdftitle={English CV Template},pdfauthor={Candidate Name},pdfsubject={Curriculum Vitae}}
\pagestyle{fancy}
\fancyhf{}
\renewcommand{\headrulewidth}{0pt}
\renewcommand{\footrulewidth}{0pt}
\fancyfoot[C]{\color{muted}\scriptsize SURNAME, Given Name \enspace | \enspace Target Function \enspace | \enspace \thepage/2}
\setlength{\parindent}{0pt}
\setlength{\tabcolsep}{0pt}
\setlength{\parskip}{0pt}
\linespread{1.03}
\setlength{\JustifyingParindent}{0pt}
\justifying
\setlength{\emergencystretch}{3em}
\tolerance=2500
\hyphenpenalty=8000
\exhyphenpenalty=8000
\color{textgrey}
\titleformat{\section}{\color{accent}\bfseries\large\raggedright}{}{0pt}{}[\color{accent}{\titlerule[1.5pt]}]
\titlespacing*{\section}{0pt}{9pt}{5pt}
\newcommand{\cvheading}[4]{%
  \par\addvspace{5pt}\noindent\begin{tabularx}{\linewidth}{@{}>{\raggedright\arraybackslash}Xr@{}}
    \textbf{#1} & \color{muted}\small #2 \\
    \textit{#3} & \color{muted}\small #4
  \end{tabularx}\par\vspace{-1pt}}
\newcommand{\cvheadingfirst}[4]{%
  \par\noindent\begin{tabularx}{\linewidth}{@{}>{\raggedright\arraybackslash}Xr@{}}
    \textbf{#1} & \color{muted}\small #2 \\
    \textit{#3} & \color{muted}\small #4
  \end{tabularx}\par\vspace{-1pt}}
\newenvironment{cvitems}{%
  \begin{itemize}[leftmargin=4.6mm,label=\raisebox{0.25ex}{\tiny$\bullet$},itemsep=2pt,topsep=2.2pt,parsep=0pt,partopsep=0pt]\justifying
}{\end{itemize}\vspace{1.5pt}}
\newcommand{\skillline}[2]{\justifying\textbf{#1:} #2\par\vspace{1.5pt}}

\begin{document}
{\fontsize{27}{30}\selectfont\bfseries SURNAME, Given Name}\par
\vspace{8pt}
{\color{accent}\bfseries\fontsize{12}{15}\selectfont TARGET FUNCTION \enspace | \enspace CORE EXPERTISE \enspace | \enspace INDUSTRY FOCUS}\par
\vspace{5pt}
{\small Tel: XXXXXXXXX \enspace | \enspace
\href{mailto:name@example.com}{name@example.com}\par
\href{https://www.linkedin.com/in/username}{linkedin.com/in/username} \enspace | \enspace
\href{https://github.com/username}{github.com/username}}
\par\vspace{9pt}
\justifying Write a concise two- or three-sentence profile covering scientific background, relevant platforms, evidence of delivery, and the target role. Prioritise keywords that genuinely match the vacancy.

\section{EDUCATION}
\cvheadingfirst{University Name}{20XX--Present}{PhD Candidate in Subject, Department or Institute}{ }
\begin{cvitems}
  \item Research focus: define the core question, application area, or technology platform in one line.
  \item Thesis submission expected Month 20XX; viva expected Month 20XX.
\end{cvitems}
\cvheading{University Name}{20XX--20XX}{MSc in Subject; final classification}{ }
\begin{cvitems}
  \item Relevant training: list three to five modules aligned with the target role.
  \item State a dissertation distinction, award, or other evidence of academic performance separately.
\end{cvitems}
\cvheading{University Name}{20XX--20XX}{BEng/BSc in Subject; GPA X.X/X.X}{ }
\begin{cvitems}
  \item Relevant training: concise selection of technical, analytical, or regulatory modules.
  \item Include scholarships, thesis awards, or competitions only when they strengthen the application.
\end{cvitems}

\section{CORE RESEARCH EXPERIENCE}
\cvheadingfirst{Research Project Title}{20XX--Present}{Doctoral Researcher, University or Research Institute}{ }
\begin{cvitems}
  \item Investigate a defined question, process, or mechanism using appropriate models and controls.
  \item Design experiments across variables, conditions, time points, and repeated measurements; own data quality and interpretation.
  \item Apply method A, method B, and method C; clarify the work performed independently and any cross-functional contribution.
  \item Quantify the strongest outcome using effect size, throughput, reproducibility, accuracy, or a decision enabled by the work.
\end{cvitems}
\cvheading{Research Project Title}{20XX.XX--20XX.XX}{Research Assistant, Organisation Name}{ }
\begin{cvitems}
  \item Validated candidate targets or optimised a workflow using technique A, technique B, and technique C.
  \item Built dose/time-course studies and translated results into a repeatable experimental plan and stakeholder update.
\end{cvitems}
\cvheading{Dissertation Project Title}{20XX.XX--20XX.XX}{MSc Research Project, University $\times$ Partner Institute}{Project grade/award}
\begin{cvitems}
  \item Established or applied a quantitative assay, imaging workflow, or analysis pipeline across defined samples and metrics.
  \item Report the most persuasive result and explain its scientific or translational relevance.
\end{cvitems}

\newpage
\begin{flushright}
  {\Large\bfseries SURNAME, Given Name}\par
  \vspace{2pt}
  {\color{muted}\small name@example.com \enspace | \enspace Tel: XXXXXXXXX}
\end{flushright}
\vspace{-10pt}

\section{ADDITIONAL RESEARCH \& INDUSTRY EXPERIENCE}
\cvheadingfirst{Research or Innovation Project}{20XX.XX--20XX.XX}{Project Lead/Team Member, Organisation Name}{ }
\begin{cvitems}
  \item Led or contributed within a team of X to deliver a defined scientific, process, or platform objective.
  \item Quantify an assay established, workflow improved, candidate screened, or project milestone delivered.
\end{cvitems}
\cvheading{Applied, Process, or Quality Project}{20XX.XX--20XX.XX}{Research Project, University $\times$ Company}{Project award}
\begin{cvitems}
  \item Summarise option screening, method development, process optimisation, or quality activities.
  \item Name the key instruments and acceptance criteria, followed by the decision or recommendation produced.
\end{cvitems}
\cvheading{Official Company Name}{20XX.XX--20XX.XX}{Quality Control / R\&D Intern}{ }
\begin{cvitems}
  \item Supported sample testing, stability studies, documentation, or quality-system activities under supervision.
  \item Specify the method, sample type, and personal responsibility; avoid generic claims of familiarity.
\end{cvitems}
\cvheading{Official Company Name}{20XX.XX--20XX.XX}{Research Laboratory Assistant / Intern}{ }
\begin{cvitems}
  \item Supported experimental execution, process optimisation, or analytical testing of development samples.
  \item Where discloseable, quantify throughput, efficiency, consistency, or turnaround improvement.
\end{cvitems}

\section{TECHNICAL SKILLS}
\skillline{Core Methods}{Method A; method B; method C; method D; method E.}
\skillline{Technical Platforms}{Platform A; platform B; platform C; instrumentation or systems.}
\skillline{Data Analysis}{Data processing; visualisation; statistics; programming language.}
\skillline{Software \& Tools}{Professional software; project tools; quality systems; compliance fundamentals.}

\section{LEADERSHIP \& ORGANISATIONAL EXPERIENCE}
\cvheadingfirst{Student Society or Professional Association}{20XX.XX--20XX.XX}{Executive Committee Member / Department Lead}{ }
\begin{cvitems}
  \item Led an X-person team, coordinating priorities, delegation, and delivery across internal and external stakeholders.
  \item Delivered X events for X attendees while managing budget, communications, suppliers, or operational risk.
\end{cvitems}
\cvheading{Society or Volunteer Organisation}{20XX.XX--20XX.XX}{President / Project Lead}{ }
\begin{cvitems}
  \item Describe the organisation's scale and the operating process or governance improvement you introduced.
  \item Highlight transferable leadership, negotiation, communication, and problem-solving evidence.
\end{cvitems}

\section{LANGUAGES \& ADDITIONAL INFORMATION}
\skillline{Languages}{Language A: native \enspace | \enspace Language B: professional proficiency (test score if useful).}
\skillline{Additional}{Selected awards, certifications, professional memberships, or role-relevant information.}
\end{document}
